% Clase 22 --- El desfasaje en función de la frecuencia (§4.2).
% El argumento de Z pasa de -pi/2 a +pi/2 y vale CERO justo en resonancia. Con
% R chica la transición se vuelve casi discontinua: el circuito cambia de
% capacitivo a inductivo en un suspiro.
\begin{center}
\begin{tikzpicture}[scale=1]

  \begin{axis}[
    fisfig, width=10.5cm, height=5.2cm,
    xlabel={$\omega$}, ylabel={$\varphi = \arg\hat Z$},
    xmin=0.05, xmax=2.6, ymin=-1.85, ymax=1.85,
    xtick={1}, xticklabels={$\frac{1}{\sqrt{LC}}$},
    ytick={-1.5708,0,1.5708},
    yticklabels={$-\frac{\pi}{2}$,$0$,$\frac{\pi}{2}$},
    legend style={at={(0.97,0.28)}, anchor=east, draw=figgray!40,
                  fill=white, font=\footnotesize},
  ]
    \addplot[figred, smooth, samples=300, domain=0.05:2.6]
      {rad(atan((x - 1/x)/0.2))};
    \addlegendentry{$R$ chica}
    \addplot[figblue, smooth, samples=300, domain=0.05:2.6]
      {rad(atan((x - 1/x)/1.1))};
    \addlegendentry{$R$ grande}
    \draw[figaux] (axis cs:1,-1.85) -- (axis cs:1,1.85);
    \draw[figaux] (axis cs:0.05,0) -- (axis cs:2.6,0);
    \node[figlblaux, anchor=south west] at (axis cs:1.5,1.0)
      {inductivo};
    \node[figlblaux, anchor=north east] at (axis cs:0.65,-1.0)
      {capacitivo};
  \end{axis}

  \node[figlbl, anchor=north, align=center] at (5.2,-1.6)
    {$\varphi = \arctan\dfrac{\omega L - \dfrac{1}{\omega C}}{R}$,\qquad
     $I(t) = I_M\cos(\omega t - \varphi)$};

\end{tikzpicture}\\[0.5em]
{\small\itshape El argumento de la impedancia completa el cuadro que empezó el
módulo. A frecuencias bajas manda el condensador, la reactancia es negativa y
$\varphi \to -\pi/2$: la corriente adelanta a la tensión casi un cuarto de
período. A frecuencias altas manda la bobina y pasa lo contrario,
$\varphi \to +\pi/2$. Y en el medio, exactamente en resonancia, los dos efectos
se cancelan, $\varphi = 0$ y el circuito se comporta ---en cuanto a fase--- como
si sólo hubiera una resistencia: corriente y tensión van juntas, que es también
la condición para que la potencia entregada sea máxima. La resistencia controla
cuán brusco es el cambio: con $R$ grande la transición es suave y con $R$ chica
se vuelve casi un salto, hasta el punto de que en el límite $R\to 0$ el
desfasaje pasa de $-\pi/2$ a $+\pi/2$ de manera discontinua ---que es la versión
en fases de la asíntota infinita de la primera figura---. Con esto el circuito
de corriente alterna queda descrito por completo: el módulo de $\hat Z$ da la
amplitud y su argumento el desfasaje.}
\end{center}
